\section{Electroweak embedding as a projective datum}
\label{sec:ew-embedding}

The gauge-Higgs sector gives the tree-level vector masses
\begin{equation}
  M_W^2=\frac{g^2v^2}{4},
  \qquad
  M_Z^2=\frac{(g^2+g'^2)v^2}{4},
  \qquad
  M_\gamma^2=0.
  \label{eq:tree-masses}
\end{equation}
The ratio is
\begin{equation}
  1-\frac{M_W^2}{M_Z^2}=\frac{g'^2}{g^2+g'^2}.
  \label{eq:tree-angle}
\end{equation}
Thus the vacuum scale \(v\) fixes the radial size of the massive vector pair, while the quotient fixes a projective direction in the two-coupling plane.

The DeVries proposal is read in the same projective language:
\begin{equation}
  \frac{M_W^2}{M_Z^2}\quad \longleftrightarrow\quad \frac{x_+(3/4)}{x_+(2)}.
  \label{eq:projective-identification}
\end{equation}
The common scale is then supplied by the electroweak vacuum or, in a deeper model, by the compactification or brane scale that generates the effective gauge-Higgs sector.

The unbroken limit is the ray
\begin{equation}
  v\to0
  \quad\text{with}\quad
  \frac{g'^2}{g^2+g'^2}\;\text{fixed}.
  \label{eq:ew-ray}
\end{equation}
Along this ray, \(M_W\) and \(M_Z\) vanish together and the photon remains massless.  This is the electroweak deformation relevant to the construction.  The projective datum remains meaningful as the limiting direction even when the radial mass scale is taken to zero.

The central analytical question is the ordered assignment
\begin{equation}
  (J_W,J_Z)=\left(\frac34,2\right).
  \label{eq:Jassignment}
\end{equation}
The working interpretation is that these entries encode electroweak representation data associated with the symmetry-breaking direction.  The \(J=3/4\) datum is naturally associated with a doublet-like \(T=1/2\) object, while \(J=2\) is naturally associated with an adjoint-like \(T=1\) object.  This statement is a candidate route rather than a completed derivation.  A completed derivation must start from the gauge-Higgs Lagrangian or from a compactification boundary problem and produce Eq.~\eqref{eq:Jassignment} without fitting to the observed masses.

\paragraph*{Section status.}  The projective interpretation follows from the electroweak mass matrix.  The representation-theoretic origin of Eq.~\eqref{eq:Jassignment} remains the primary open calculation.
